% =============================================================================
%  s03_conclusiones.tex — Sección 3: Conclusiones y recomendaciones
% =============================================================================

\section{Conclusiones}

% ── Frame: Recomendaciones ───────────────────────────────────────────────────
% NOTA \pause: descomenta para revelar items en vivo. En PDF genera páginas duplicadas.
\begin{frame}{Puntos de Conclusion}
  \begin{enumerate}
    \item \textbf{Eje Social (Identidad, Resistencia y la Reivindicación del Poblador):} \\
          Los humedales representan un nodo sociocultural donde la identidad se debate entre la preservación y la exclusión. 
          El sistema de chinampas en Xochimilco es el ejemplo más claro de una interacción equilibrada entre la especie humana 
          y el ambiente, favoreciendo el crecimiento de especies endémicas \cite{paot2006xochimilco}.
    \item \textbf{Eje de Planeación (El Fracaso de la Expansión Horizontal y el Acceso a la Ciudad): } \\
          El crecimiento de la Zona Metropolitana del Valle de México hacia sus periferias revela una incapacidad de planeación vertical, 
          forzando a las zonas limítrofes a convertirse en áreas de exclusión. Mientras que en Xochimilco la infraestructura gris como el puente de 
          Cuemanco amenaza con fragmentar el territorio \cite{segob2020proposicion}, 
          en Aragón la urbanización temprana (1970-1990) consumió gran parte de su identidad original.
  \end{enumerate}
\end{frame}

\begin{frame}{Puntos de Conclusion}
  \begin{enumerate}
    \item \textbf{Eje de Movilidad (El Transporte como Ente Regulador y no Destructor):} \\
          La palabra "progreso" se ha vinculado históricamente con la destrucción de ecosistemas, pero bajo una visión de sustentabilidad, 
          debe transitar hacia la regulación. El transporte masivo es una necesidad básica cuya ineficiencia afecta la escala macro de la metrópoli. 
          Bajo el \textit{Modelo VFT}, la propuesta de un corredor circular sobre el Anillo Periférico debe actuar como una frontera estructural \cite{garibaldi2023vft}. 
    \item Al gestionar matemáticamente la fuerza capilar del transporte y consolidar los flujos en nodos optimizados, la arquitectura vial puede fungir como una barrera 
          artificial que detenga la dispersión de transeúntes hacia áreas naturales protegidas \cite{garibaldi2023vft}.
  \end{enumerate}
\end{frame}

% ── Frame: Cierre ────────────────────────────────────────────────────────────
\begin{frame}[plain]
  \begin{center}
    \vfill
    {\color{ColorAcento}\rule{0.6\textwidth}{1.5pt}}\\[0.5cm]
    {\Large\bfseries\color{ColorPrincipal} Gracias}\\[0.3cm]
    {\normalsize \Autor}\\[0.1cm]
    {\small \MateriaCorta}\\[0.1cm]
    {\small \Semestre}\\[0.5cm]
    {\color{ColorAcento}\rule{0.6\textwidth}{1.5pt}}
    \vfill
  \end{center}
\end{frame}
